\documentclass[11pt,letterpaper]{article}

% Self-contained source: no external bibliography, images, or custom class.
\usepackage[T1]{fontenc}
\usepackage{lmodern}
\usepackage[margin=1in]{geometry}
\usepackage{amsmath,amssymb,amsthm,mathtools}
\usepackage{booktabs}
\usepackage{enumitem}
\usepackage{microtype}
\usepackage{xcolor}
\usepackage{hyperref}

\definecolor{linkblue}{RGB}{28,65,110}
\hypersetup{
  colorlinks=true,
  linkcolor=linkblue,
  citecolor=linkblue,
  urlcolor=linkblue,
  pdftitle={Hidden Blocking Pairs: Auditing Intermediaries in Housing Markets},
  pdfsubject={Research note and proof sketch for Harvard CS 2880},
  pdfkeywords={stable matching, intermediary, housing, disclosure, NP-completeness}
}
\urlstyle{same}
\setlength{\emergencystretch}{2em}
\setlist[itemize]{leftmargin=1.5em,itemsep=2pt,topsep=4pt}
\setlist[enumerate]{leftmargin=1.8em,itemsep=3pt,topsep=4pt}
\allowdisplaybreaks[1]

\newtheorem{theorem}{Theorem}[section]
\newtheorem{proposition}[theorem]{Proposition}
\newtheorem{corollary}[theorem]{Corollary}
\theoremstyle{definition}
\newtheorem{definition}[theorem]{Definition}
\newtheorem{example}[theorem]{Example}
\theoremstyle{remark}
\newtheorem{remark}[theorem]{Remark}

\newcommand{\DA}{\operatorname{DA}_{B}}
\newcommand{\BP}{\operatorname{BP}}
\newcommand{\SD}{\textsc{StableDisclosure}}
\newcommand{\NP}{\mathsf{NP}}
\newcommand{\botv}{\bot}
\newcommand{\Feas}{\mathcal{F}}
\newcommand{\R}{\mathbb{R}}

\title{\vspace{-1.5em}\textbf{Hidden Blocking Pairs}\\[4pt]
  \large Auditing Intermediaries in Housing Markets}
\author{\normalsize Harvard CS 2880 \quad Research note and proof sketch}
\date{\normalsize October 25, 2026}

\begin{document}
\maketitle
\vspace{-1.4em}

\begin{abstract}
Housing intermediaries may influence which feasible transactions their clients
discover. This creates a distinction between stability in the disclosed market
and stability under clients' full opportunities. We formalize a minimum-cost
disclosure problem: audit brokerage firms, reveal the transactions incident to
their clients, and require buyer-proposing deferred acceptance to return a
matching stable in the full market. A reduction from bounded-degree
\textsc{Vertex Cover} establishes NP-completeness even when every feasible-market
component is a four-vertex path, preference lists have length at most two, and
each firm represents at most three clients. As a complementary positive result,
a shared binary partnership-bonus model reduces simultaneous stability to
super-stable matching with partial orders. These observations provide a starting
point for studying the computational and distributional consequences of
intermediary-controlled information.
\end{abstract}

\noindent\textbf{Scope and status.}
The statements proved below concern the explicit models in this note.
They do not establish novelty relative to the matching, information-acquisition,
or brokerage literatures. The proposed research contribution is an intervention
problem and its tractability boundaries; the presence of four economic roles,
by itself, is not a hardness result.

\section{Motivation and modeling choices}

There are four economic roles: a buyer, the buyer's agent, a seller's agent, and
a seller. An intermediary may favor a transaction because of its commission,
expected closing effort, or a relationship with the other intermediary. Such
relationships can also lower transaction costs. The model therefore does not
assume that professional collaboration is inherently harmful.

Commission-driven steering has an empirical precedent in residential brokerage
\cite{barwick2017}. Our narrower theoretical question is:
\begin{quote}
\emph{When intermediaries limit access to feasible transactions, what disclosure
intervention makes the resulting allocation stable for their clients?}
\end{quote}

We initially study purchases, with one property per seller, fixed representation,
and fixed transaction terms. Prices, commissions, financing requirements, and
closing conditions enter feasibility or client preferences. Bargaining,
endogenous representation, and broker-wide complementarities are outside this
initial model. A rental application would require a separate interpretation of
tenants, landlords, and intermediary incentives.

The classical stable-marriage problem is polynomial-time solvable:
deferred acceptance makes at most one proposal per acceptable pair
\cite{gale1962}. Thus a complexity claim must identify an additional decision or
optimization problem.

\section{Client stability and disclosure}

Let \(B\) and \(S\) be finite sets of buyers and sellers. Let
\(E\subseteq B\times S\) be an explicitly represented graph of mutually
acceptable transactions. Each client \(i\in B\cup S\) has a strict preference
order \(\succ_i\) over its acceptable partners, followed by the outside option
\(\botv\). In particular, every edge in \(E\) is preferred to remaining unmatched.

A matching \(M\subseteq E\) assigns at most one partner to each client.
Write \(M(i)=\botv\) if \(i\) is unmatched. For any \(H\subseteq E\), define
\begin{equation}
\BP_H(M)=
\left\{(b,s)\in H\setminus M:
s\succ_b M(b)\ \text{and}\ b\succ_s M(s)
\right\}.
\label{eq:blocking}
\end{equation}
We use this definition when \(M\) is a matching in \(H\), or in a subset of \(H\).

\begin{definition}[Visible and full-market stability]
Let \(E_0\subseteq E\) be the initially disclosed graph.
A matching \(M\subseteq E_0\) is \emph{visibly stable} if
\(\BP_{E_0}(M)=\varnothing\), and \emph{client-stable in the full market} if
\(\BP_E(M)=\varnothing\).
\end{definition}

For example, suppose \(E_0=\{b_1s_1,b_2s_2\}\), while the full market also
contains \(b_1s_2\) and \(b_2s_1\). If every client prefers its partner in the
crossed arrangement, the disclosed matching is visibly stable but admits two
full-market blocking pairs.

\begin{proposition}[Unanimous four-party consent is not a hardness source]
Suppose representation is fixed and a blocking transaction requires both
clients and both intermediaries to strictly prefer that transaction.
Every client-stable matching has no such unanimous blocking transaction.
Consequently, an outcome without a unanimous block can be found in polynomial
time whenever these are the only deviation constraints.
\end{proposition}

\begin{proof}
A unanimous blocking transaction would, in particular, be a blocking pair for
the buyer and seller. A client-stable matching rules out that pair.
Run deferred acceptance on \(E\) to obtain such a matching.
\end{proof}

This observation separates a four-party consent requirement from the
information restriction studied below.

\subsection{A firm-level disclosure intervention}

Let \(\Feas\) be a set of brokerage firms. Fixed maps
\[
f_B:B\longrightarrow\Feas,
\qquad
f_S:S\longrightarrow\Feas
\]
record the firms employing the clients' agents. The buying and listing agents
remain distinct economic roles; a single firm may represent clients in both
roles. Define a firm's portfolio by
\[
C_f=\{b\in B:f_B(b)=f\}\cup\{s\in S:f_S(s)=f\}.
\]

An audit of firm \(f\) makes all feasible transactions incident to \(C_f\)
available to their endpoints. For audited firms \(A\subseteq\Feas\), set
\begin{equation}
E(A)=E_0\cup
\left\{(b,s)\in E:
f_B(b)\in A\ \text{or}\ f_S(s)\in A
\right\}.
\label{eq:disclosure}
\end{equation}
The disjunction in \eqref{eq:disclosure} is an assumption: auditing either
endpoint's firm suffices to expose the opportunity. An audit expands access; it
does not force either client to accept a transaction.

Let
\begin{equation}
M(A)=\DA\bigl(E(A),\succ|_{E(A)}\bigr)
\label{eq:response}
\end{equation}
be buyer-proposing deferred acceptance using true client preferences restricted
to the disclosed graph. Fix any deterministic implementation; with strict
preferences, the buyer-optimal stable outcome is well-defined.

\begin{definition}[\SD]
An instance consists of \(B,S,E,E_0,\succ,\Feas,f_B,f_S\), positive integer
audit costs \((c_f)_{f\in\Feas}\), and a nonnegative integer budget \(K\).
All graphs and preference lists are explicit, and integers are binary encoded.
The question is whether there is \(A\subseteq\Feas\) such that
\[
\sum_{f\in A}c_f\leq K
\qquad\text{and}\qquad
\BP_E(M(A))=\varnothing.
\]
The optimization version minimizes \(\sum_{f\in A}c_f\) subject to the same
stability requirement.
\end{definition}

The designer is given \(E\) and the full preference profile. This is an
\emph{offline intervention benchmark}; it does not assume that clients already
know those opportunities. An adaptive investigator who learns edges or
preferences through audits faces a different information-acquisition problem.

\section{NP-completeness of selecting disclosures}

\begin{theorem}
\label{thm:hardness}
\SD{} is NP-complete. This remains true under all of the following restrictions:
\begin{enumerate}[label=(\roman*)]
\item every audit has unit cost;
\item every firm represents at most three clients in total;
\item every connected component of the full feasible graph \(E\) is a
      four-vertex path;
\item every client has at most two acceptable partners; and
\item \(E_0\) is a perfect matching of all clients.
\end{enumerate}
The matching produced after disclosure is allowed to leave clients unmatched.
\end{theorem}

\begin{proof}
\textbf{Membership in \(\NP\).}
A certificate is an audit set \(A\). Its cost can be checked in polynomial time.
Construct \(E(A)\), run deferred acceptance, and check every edge of \(E\) using
\eqref{eq:blocking}. With inverse-rank tables, the proposal process and the
blocking-pair scan each take \(O(|E|)\) time, apart from input preprocessing.
Thus a successful audit set is polynomial-time verifiable.

\medskip
\noindent\textbf{Reduction.}
We reduce from \textsc{Vertex Cover} on simple graphs of maximum degree three,
a known NP-complete problem \cite{garey1976}. Let its instance be
\((G=(V,H),k)\). For each \(v\in V\), create a firm \(f_v\).
Orient each graph edge arbitrarily, writing \(e=(u,v)\) solely to distinguish
its two endpoint roles.

For each \(e=(u,v)\in H\), create buyers \(b_e,t_e\), sellers \(s_e,r_e\), and
exactly three feasible transactions:
\begin{equation}
E_e=\{b_es_e,\ b_er_e,\ t_es_e\},
\qquad
(E_0)_e=\{b_er_e,\ t_es_e\}.
\label{eq:gadget}
\end{equation}
Here \(b_es_e\) abbreviates the ordered buyer--seller pair \((b_e,s_e)\).
The complete acceptable preference lists are
\begin{equation}
\begin{array}{c@{\qquad}c}
b_e:\ s_e\succ r_e\succ\botv
&
t_e:\ s_e\succ\botv
\\[3pt]
s_e:\ b_e\succ t_e\succ\botv
&
r_e:\ b_e\succ\botv .
\end{array}
\label{eq:preferences}
\end{equation}
There are no feasible transactions between different gadgets.

Assign \(b_e\) to firm \(f_u\), and \(s_e\) to firm \(f_v\).
Assign \(t_e\) and \(r_e\) to two fresh, distinct auxiliary firms, used in no
other gadget. Set every audit cost to one, and set the disclosure budget to
\(K=k\). Auditing an auxiliary firm reveals no new edge, because its client's
only feasible transaction is already in \(E_0\).

\medskip
\noindent\textbf{Local behavior.}
For any audit set \(A\), the central transaction \(b_es_e\) is exposed if and
only if \(f_u\in A\) or \(f_v\in A\). The two possible responses are:
\begin{center}
\begin{tabular}{@{}p{0.28\linewidth}p{0.29\linewidth}p{0.32\linewidth}@{}}
\toprule
Audits at the endpoints & Deferred-acceptance outcome & Full-market stability\\
\midrule
Neither \(f_u\) nor \(f_v\)
&
\(\{b_er_e,t_es_e\}\)
&
Blocked by \(b_es_e\).
\\[5pt]
At least one of \(f_u,f_v\)
&
\(\{b_es_e\}\); \(t_e,r_e\) unmatched
&
Stable in \(E_e\).
\\
\bottomrule
\end{tabular}
\end{center}
In the second row, \(b_e\) and \(s_e\) are mutual first choices. The other
feasible pairs cannot block: \(b_e\) prefers \(s_e\) to \(r_e\), and \(s_e\)
prefers \(b_e\) to \(t_e\). Since gadgets are disconnected, their
deferred-acceptance outcomes and blocking conditions are independent.

\medskip
\noindent\textbf{Correctness.}
If \(C\subseteq V\) is a vertex cover with \(|C|\leq k\), audit
\(A_C=\{f_v:v\in C\}\). Every graph edge has an audited endpoint, so every gadget
has the stable outcome in the second row. Hence \(M(A_C)\) is stable in \(E\)
and its audit cost is at most \(K\).

Conversely, suppose an audit set \(A\) of cost at most \(K\) produces an outcome
stable in \(E\). Remove all auxiliary firms from \(A\); this changes neither
disclosures nor the outcome. For each graph edge \(e=(u,v)\), stability requires
\(f_u\in A\) or \(f_v\in A\), since otherwise \(b_es_e\) blocks.
Therefore
\[
C_A=\{v\in V:f_v\in A\}
\]
is a vertex cover of size at most \(k\).

\medskip
\noindent\textbf{Size and restrictions.}
The construction has \(4|H|\) clients, \(3|H|\) feasible transactions, and
\(|V|+2|H|\) firms, and can be constructed in polynomial time. A vertex firm
represents one client per incident graph edge, hence at most three clients.
An auxiliary firm represents one client. Each feasible component is the path
\(r_e-b_e-s_e-t_e\); every preference list has length at most two.
The two outer edges in every gadget form a perfect matching in \(E_0\).
All claimed restrictions follow.
\end{proof}

\begin{corollary}
The minimum-cost disclosure optimization problem is NP-hard, even under the
restrictions of Theorem~\ref{thm:hardness}.
\end{corollary}

\begin{remark}[What the reduction establishes]
The hard task is selecting firm-level disclosures before a polynomial-time
matching procedure. It is not computing an ordinary stable matching. The
reduction uses a single audit action that may expose transactions for buyers
and sellers represented by the same firm. It does not establish hardness when
buying-agent and listing-agent audit actions are required to be disjoint.
\end{remark}

\begin{remark}[Stability and participation are separate objectives]
Initially every gadget has two transactions, while its full-market stable
outcome has only one. Thus this proof does not show that disclosure increases
transaction volume, improves every client's outcome, or increases welfare.
A requirement that the final matching be perfect would define a different
problem, to which this reduction does not apply.
\end{remark}

For this constructed family, the minimum number of audits equals the minimum
vertex-cover size. This correspondence is not an approximation algorithm for
arbitrary disclosure instances: rematching can couple transactions in ways the
independent gadgets do not represent.

\section{A tractable partnership-bonus baseline}

Consider a separate model in which all feasible transactions are available,
but intermediary relationships alter recommendation scores. Keep the same
acceptable graph \(E\). For every endpoint \(i\) of transaction \(e\), let
\(u_i(e)\in\mathbb{Q}\) be its base score. Set \(u_i(\botv)\) strictly below all its
acceptable scores.

Each transaction has a shared label \(F_e\in\{0,1\}\): both endpoints agree on
whether the intermediaries are partners. For client-specific constants
\(\lambda_i\in\mathbb{Q}_{\geq0}\), define
\begin{equation}
\widetilde u_i(e)=u_i(e)+\lambda_iF_e.
\label{eq:bonus}
\end{equation}
The outside-option score is unchanged, and \(F_{\botv}=0\).
All rational inputs have binary encodings.
Let \(P_i\) and \(Q_i\) be the rankings induced by \(u_i\) and
\(\widetilde u_i\), respectively. Assume both profiles are strict; for this
proposition, score equalities at these two endpoints of the parameter interval
are excluded.

A ranking in \(Q\) is a mediated recommendation ranking, not automatically a
client's true preference or a broker's portfolio-wide utility.
Requiring stability under both \(P\) and \(Q\) is a deliberate robustness
criterion, distinct from unanimous four-party consent.

Define an intersection relation on each client's acceptable partners and
outside option:
\begin{equation}
x\succ_i^* y
\quad\Longleftrightarrow\quad
x\succ_i^P y\ \text{and}\ x\succ_i^Q y.
\label{eq:intersection}
\end{equation}
This relation is a strict partial order.

\begin{definition}[Super-stability for the intersection orders]
A matching \(M\subseteq E\) is \emph{super-stable} under \(\succ^*\) if every
unmatched feasible pair \((b,s)\in E\setminus M\) has an endpoint that strictly
prefers its current assignment in the partial order:
\begin{equation}
M(b)\succ_b^* s
\quad\text{or}\quad
M(s)\succ_s^* b.
\label{eq:superstable}
\end{equation}
An incomparable alternative is therefore treated as a possible reason to
deviate, rather than as a certified rejection.
\end{definition}

\begin{theorem}[Binary bonuses reduce to partial-order super-stability]
\label{thm:binary}
A matching is stable under both \(P\) and \(Q\) if and only if it is
super-stable under the intersection orders \(\succ^*\).
Consequently, existence and construction of a simultaneously stable matching
are polynomial-time solvable.
\end{theorem}

\begin{proof}
For a matching \(M\), let \(e_i^M\) denote the transaction containing client \(i\),
or \(\botv\) if unmatched. For a candidate transaction \(e\in E\setminus M\),
define the endpoint's gain under each profile:
\begin{align}
\Delta_i^P(e)
 &=u_i(e)-u_i(e_i^M),\\
\Delta_i^Q(e)
 &=\Delta_i^P(e)+\lambda_i\bigl(F_e-F_{e_i^M}\bigr).
\label{eq:gain}
\end{align}

Suppose first that \(M\) is stable under \(P\) and \(Q\).

If \(F_e=1\), equation~\eqref{eq:gain} gives
\(\Delta_i^Q(e)\geq\Delta_i^P(e)\) at both endpoints, since
\(F_{e_i^M}\in\{0,1\}\). Stability under \(Q\), together with strict
preferences, supplies an endpoint \(i\) with \(\Delta_i^Q(e)<0\).
Then \(\Delta_i^P(e)\leq\Delta_i^Q(e)<0\), so that same endpoint rejects
\(e\) under both profiles.

If \(F_e=0\), then \(\Delta_i^Q(e)\leq\Delta_i^P(e)\) at both endpoints.
Use stability under \(P\) to find an endpoint with
\(\Delta_i^P(e)<0\). That endpoint also satisfies \(\Delta_i^Q(e)<0\),
and again rejects \(e\) under both profiles.

Thus every unmatched feasible pair has a common rejecting endpoint, exactly
as required by \eqref{eq:superstable}.

Conversely, if \eqref{eq:superstable} holds, its rejecting endpoint rejects
the candidate pair under \(P\) and under \(Q\). Hence no pair blocks in either
profile.

Finally, the orders \(\succ_i^*\) can be constructed in polynomial time by
intersecting two strict rankings. Super-stable matching with partial orders
on both sides is polynomial-time solvable; the algorithmic results of
Cseh and Juhos cover this case \cite{cseh2021}. Applying that algorithm gives
a matching or certifies that none exists.
\end{proof}

\begin{corollary}[Robustness to bounded incentive strengths]
\label{cor:robust}
Under the assumptions of Theorem~\ref{thm:binary}, simultaneous stability under
\(P,Q\) is equivalent to having no strictly blocking pair for every vector
\[
(\theta_i)_{i\in B\cup S}\in
\prod_{i\in B\cup S}[0,\lambda_i],
\qquad
u_i^\theta(e)=u_i(e)+\theta_iF_e.
\]
Intermediate score ties are allowed; a blocking pair must strictly improve
both clients' scores.
\end{corollary}

\begin{proof}
Theorem~\ref{thm:binary} gives each candidate pair an endpoint that rejects it
strictly at \(\theta_i=0\) and \(\theta_i=\lambda_i\).
The score difference between its current match and the candidate is affine
in \(\theta_i\), so it remains strictly positive throughout the interval.
The pair therefore cannot block for any vector \(\theta\).
The converse follows by choosing the vectors \(\theta=0\) and
\(\theta=(\lambda_i)_i\).
\end{proof}

The theorem is an existence test, not an existence guarantee. For example, in
a complete \(2\times2\) market, let every client's base favorite be its diagonal
partner, with utility one, and the other partner have utility zero. Label the
two crossed transactions \(F=1\), the diagonal transactions \(F=0\), and set
all bonuses to two. Then \(P\) has the unique diagonal stable matching and
\(Q\) the unique crossed stable matching. There is no simultaneous solution.

\subsection{What changes with three bonus levels?}

The common-rejecting-endpoint proof can fail if bonuses take values
\(0,1,2\). Here is an explicit example. Let
\[
B=\{a,t\},\quad S=\{r,s\},\quad
E=\{ar,as,ts\},\quad M=\{ar,ts\},
\]
and set
\[
F_{ar}=2,\qquad F_{as}=1,\qquad F_{ts}=0,\qquad
\lambda_i=2\ \text{for every client}.
\]
Choose
\[
u_a(r)=0,\quad u_a(s)=1,\qquad
u_s(a)=0,\quad u_s(t)=1,
\]
with \(u_r(a)=u_t(s)=0\) and outside-option utility \(-1\).

The only candidate pair is \(as\). Under \(P\), seller \(s\) rejects it.
Under \(Q\), buyer \(a\) rejects it because
\(\widetilde u_a(ar)=4>3=\widetilde u_a(as)\).
Thus \(M\) is stable under both profiles, but neither endpoint rejects \(as\)
under their intersection orders. The binary reduction is therefore invalid
for this instance.

This counterexample proves failure of a particular reduction, not NP-hardness.
It also does not imply hardness of the stronger independent-interval
robustness problem. For arbitrary affine score perturbations with independent
endpoint parameters, rejection throughout an interval is still captured by
the intersection of its endpoint rankings.

\section{Relation to existing work and the research target}

Several established models delimit the contribution.

\begin{itemize}
\item \emph{Information and coordination.}
Socially stable matching distinguishes communication-enabled blocking pairs
from other potential deviations \cite{askalidis2013}. Our disclosure model
also restricts available transactions and selects grouped disclosure actions
before a specified matching procedure.
\item \emph{Brokerage incentives.}
Broker fees, matching choices, and instability have been studied in economic
models of brokerage \cite{boerner2023}. A new contribution must go beyond the
observation that intermediaries can have misaligned incentives.
\item \emph{Contracts and networks.}
Supply-chain matching with bilateral contracts admits stable outcomes under
appropriate substitutability conditions \cite{hatfield2012}. Adding commissions
or additional transaction layers does not automatically establish hardness.
\item \emph{Multiple preference profiles.}
Joint stability for two arbitrary profiles is already NP-complete
\cite{miyazaki2017}. Theorem~\ref{thm:binary} exploits a particular shared-bonus
structure, so arbitrary-profile hardness does not transfer automatically.
\item \emph{Multiple matching dimensions.}
Stable existence for cyclic \(k\)-dimensional preferences is NP-complete for
\(k\geq3\) \cite{lam2019}. That preference structure should not be imposed on
housing solely to inherit a known hardness result.
\end{itemize}

The recommended next step is to develop the disclosure model beyond the
independent gadgets of Theorem~\ref{thm:hardness}. Three concrete questions are:

\begin{enumerate}
\item \textbf{Find a tractability boundary.}
What happens when each firm represents clients on only one side, portfolios
have size at most two, or the graph coupling firms has restricted structure?
The proof above does not settle these cases.
\item \textbf{Design an intervention algorithm.}
Seek an approximation or parameterized algorithm for general disclosure
instances, with an explicit guarantee. Selecting a firm that exposes a
current blocking pair is a heuristic until rematching effects are analyzed.
\item \textbf{Study discovery and client outcomes.}
For an adaptive model, specify what an audit reveals and what the investigator
knows beforehand. Evaluate audit cost, full-market blocking pairs, unmatched
clients, and client utility or rank distributions. Stability, transaction
volume, welfare, and distributional impact are distinct criteria.
\end{enumerate}

A credible first deliverable is a complete model, the two proofs above, a
comparison with adjacent literature, and one substantive algorithmic or
structural result. A publication claim would additionally require a systematic
novelty check and evidence that the chosen intervention model captures a
meaningful housing-market information friction.

\begin{thebibliography}{10}
\small

\bibitem{gale1962}
D.~Gale and L.~S.~Shapley.
College admissions and the stability of marriage.
\emph{The American Mathematical Monthly}, 69(1):9--15, 1962.
\url{https://doi.org/10.1080/00029890.1962.11989827}.

\bibitem{barwick2017}
P.~J.~Barwick, P.~A.~Pathak, and M.~Wong.
Conflicts of interest and steering in residential brokerage.
\emph{American Economic Journal: Applied Economics}, 9(3):191--222, 2017.
\url{https://doi.org/10.1257/app.20160214}.

\bibitem{garey1976}
M.~R.~Garey, D.~S.~Johnson, and L.~Stockmeyer.
Some simplified NP-complete graph problems.
\emph{Theoretical Computer Science}, 1(3):237--267, 1976.
\url{https://doi.org/10.1016/0304-3975(76)90059-1}.

\bibitem{cseh2021}
\'{A}.~Cseh and A.~Juhos.
Pairwise preferences in the stable marriage problem.
\emph{ACM Transactions on Economics and Computation}, 9(1), Article~7, 2021.
\url{https://doi.org/10.1145/3434427}.
See also the author manuscript, Section~5:
\url{https://hpi.de/friedrich/docs/publications/2021/TEAC.pdf}.

\bibitem{askalidis2013}
G.~Askalidis, N.~Immorlica, and E.~Pountourakis.
Socially stable matchings.
Preprint, 2013.
\url{https://arxiv.org/abs/1302.3309}.

\bibitem{boerner2023}
L.~Boerner and D.~Quint.
Medieval matching markets.
\emph{International Economic Review}, 64(1):23--56, 2023.
\url{https://doi.org/10.1111/iere.12600}.

\bibitem{hatfield2012}
J.~W.~Hatfield and S.~D.~Kominers.
Matching in networks with bilateral contracts.
\emph{American Economic Journal: Microeconomics}, 4(1):176--208, 2012.
\url{https://doi.org/10.1257/mic.4.1.176}.

\bibitem{miyazaki2017}
S.~Miyazaki and K.~Okamoto.
Jointly stable matchings.
In \emph{28th International Symposium on Algorithms and Computation
(ISAAC 2017)}, volume~92 of LIPIcs, pages~56:1--56:12, 2017.
\url{https://doi.org/10.4230/LIPIcs.ISAAC.2017.56}.

\bibitem{lam2019}
C.-K.~Lam and C.~G.~Plaxton.
On the existence of three-dimensional stable matchings with cyclic preferences.
Preprint, 2019.
\url{https://arxiv.org/abs/1905.02844}.

\end{thebibliography}
\end{document}
